\documentclass[11pt]{article}
\usepackage{amsmath,amssymb,amsthm,mathtools}
\usepackage{bm}
\usepackage{booktabs}
\usepackage{graphicx}
\usepackage{array}
\usepackage[margin=1in]{geometry}

\newtheorem{definition}{Definition}
\newtheorem{assumption}{Assumption}
\newtheorem{proposition}{Proposition}
\newtheorem{remark}{Remark}

\DeclareMathOperator*{\argmax}{arg\,max}
\DeclareMathOperator{\diag}{diag}
\newcommand{\R}{\mathbb{R}}
\newcommand{\E}{\mathbb{E}}
\newcommand{\norm}[1]{\left\lVert #1 \right\rVert}
\newcommand{\SO}{\mathrm{SO}(3)}
\providecommand{\IEEEPARstart}[2]{#1#2}
\newcolumntype{L}[1]{>{\raggedright\arraybackslash}p{#1}}
\setlength{\emergencystretch}{3em}
\hbadness=10000
\hfuzz=50pt

\title{Estimation-Free Station-Keeping and Safe Command-Tracking\\
from Raw Range and Inertial Measurements:\\
A Unified Constrained POMDP Formulation}
\author{}
\date{}

\begin{document}
\maketitle

%============================================================
% ABSTRACT
%============================================================
\begin{abstract}
Safe quadrotor operation in GNSS-denied and perception-degraded environments
is commonly treated as collision avoidance layered onto a nominal navigation
controller, which in turn depends on an onboard pose or odometry estimate.
Such estimates are locally consistent but drift without bound in unvisited
areas and can fail outright under latency, noise, or degeneracy---precisely
the regime in which safety matters most. This paper studies a complementary
and more primitive problem: how a quadrotor can hold its local position and
safely respond to operator velocity commands without a world-frame position
setpoint, a reconstructed obstacle map, or explicit pose, odometry, or
velocity inputs to the learned policy. We formulate station-keeping as
regulation in LiDAR \emph{measurement space}: under a null command, the
vehicle anchors its current range observation as an implicit reference and
regulates subsequent measurements toward it, and the same command-conditioned
policy extends naturally to nonzero commands, tracking operator intent within
physically feasible margins. The deployed governor consumes only raw range
measurements, inertial cues, the command, and a short observation history;
its translational regulation is pose- and odometry-free, while attitude
stabilization is delegated to a conventional low-level controller. We analyze
when measurement-space station-keeping is attainable through the local range
Jacobian, showing that drift-free regulation is fundamentally limited by
geometric observability and degrades in flat-wall, tunnel, and open-field
configurations. For safety under nonzero commands, we couple the policy with
an inevitable-collision-state-inspired intervention criterion based on
clearance, approach speed, and braking feasibility. Experiments in simulated
quadrotor environments show that the method reduces station-keeping drift in
observable geometries, exposes predictable failure modes in degenerate ones,
and attenuates unsafe velocity commands while preserving tracking performance
when sensing and braking authority are sufficient.
\end{abstract}

%============================================================
% INTRODUCTION
%============================================================
\section{Introduction}

\IEEEPARstart{A}{utonomous} quadrotors are increasingly deployed in
GNSS-denied and perception-degraded environments---subterranean voids,
collapsed structures, and cluttered indoor spaces---where a human operator
issues coarse velocity intent and expects the vehicle to remain safe and
stationary between commands. The dominant design pattern treats this as a
two-stage problem: an onboard visual--inertial or LiDAR odometry front end
produces a pose estimate, a position controller closes a loop around it, and a
reactive safety filter prevents collisions. This decomposition is effective
when the estimator is healthy, but it inherits the estimator's failure modes.
Visual--inertial and LiDAR odometry are locally consistent yet accumulate
unbounded drift in previously unvisited areas\cite{sprind2025}, and onboard
localization and mapping are prone to latency, noise, and drift that can
result in estimation failure and, consequently, collision\cite{neuralcbf2024}.
In the very conditions that make these environments dangerous---featureless
walls, dust, symmetric corridors---the pose estimate is least trustworthy,
yet the safety behavior is built on top of it.

This coupling motivates a question that is seldom posed directly:
\emph{when the estimator degrades or diverges, can the vehicle still hold its
local position and safely respond to velocity commands without relying on that
estimate at all?} We argue that a useful safety layer should not fail in the
same way, and at the same time, as the localization stack it is meant to
protect. What is missing is a holding-and-response behavior that is
structurally independent of pose---one that regulates directly on what the
sensor measures rather than on a reconstructed state.

We take this independence literally. Instead of asking where the vehicle is,
we ask what it should \emph{see}. We formulate station-keeping as regulation
in LiDAR \emph{measurement space} rather than position space. Under a null
command, the vehicle freezes its current range observation as an implicit
reference and learns to maintain that range pattern; nonzero operator commands
are treated not as a separate navigation problem but as a natural extension of
the same command-conditioned policy, which induces local motion under the same
observation constraints. The deployed governor consumes only raw range
measurements, inertial cues, the command input, and a short observation
history. Its \emph{translational} regulation is pose-free and odometry-free;
we deliberately restrict this claim to translation and delegate attitude
stabilization to a conventional low-level controller, since attitude
estimation remains necessary and is not part of the learned policy's state.

Formulating holding in measurement space, however, raises a question that
position-space formulations never confront: measurement-space regulation is
only as good as the geometry that the sensor observes. Two distinct vehicle
positions can produce nearly identical range patterns, in which case no policy,
however well trained, can distinguish them or resist drift between them. We
make this precise through the local range Jacobian, which links small
translational motions to changes in LiDAR range measurements. Its conditioning
determines when station-keeping is physically attainable and predicts the
directions along which drift is unavoidable in flat-wall, tunnel, and
open-field geometries. This turns station-keeping failure from an empirical
nuisance into a predictable geometric phenomenon and, we argue, is a necessary
companion to any observation-space holding behavior.

Finally, holding is not enough: once the operator issues a nonzero command,
the governor must decide whether obeying it is safe. Because it has no map and
no global pose, this decision must also be made from local sensing. We couple
the measurement-space policy with an intervention criterion inspired by the
notion of an inevitable collision state, which evaluates whether the current
command would drive the vehicle toward an obstacle faster than it can brake
within the available clearance, accounting for both distance and approach
speed rather than distance alone. We do not claim novelty for command
filtering as such; the point is that both the holding behavior and the command
attenuation are produced from raw range and inertial sensing, without a
world-frame position estimate anywhere in the loop.

In summary, this paper makes the following contributions.

\begin{itemize}
  \item \textbf{Measurement-space station-keeping without a world-frame
  setpoint.} We cast quadrotor station-keeping and command tracking as a
  single command-conditioned regulation problem in LiDAR measurement space. A
  zero command freezes the current range pattern as an implicit reference and
  the policy maintains it from range--inertial history alone, while a nonzero
  command induces local motion under the same observation constraints---
  unifying local holding and velocity tracking without a global goal, map, or
  pose setpoint, and without any external pose, odometry, or velocity estimate
  for translational regulation.

  \item \textbf{Range-Jacobian observability analysis for pose-free holding,
  with a quantitative drift study.} We provide a local observability analysis
  based on the range Jacobian, characterizing when range-only station-keeping
  is physically attainable and exposing a fundamental limitation: in
  flat-wall, tunnel, and open-field geometries, some translational directions
  are weakly observable or unobservable, causing unavoidable drift even for an
  ideally trained policy. We validate this by relating station-keeping drift
  to the conditioning of the range Jacobian across corners, flat walls,
  corridors, and open scenes.

  \item \textbf{An ICS-inspired command-intervention layer coupled to the
  measurement-space policy.} We introduce an inevitable-collision-state-
  inspired criterion that decides, from local sensing alone, whether a
  commanded velocity remains brakeable within the available clearance,
  accounting for both obstacle distance and approach speed. Our contribution
  is not the filter in isolation but its integration with a pose-free
  measurement-space governor, so that unsafe commands are attenuated, slowed,
  or rejected before collision becomes imminent while tracking is preserved
  when sensing and braking authority are sufficient.
\end{itemize}

The remainder of this paper is organized as follows.
Section~\ref{sec:related} positions our work against learned barrier
functions, observation-space control synthesis, and safe velocity filtering.
Section~\ref{sec:formulation} formalizes measurement-space station-keeping and
the command-conditioned policy. Section~\ref{sec:observability} develops the
range-Jacobian observability analysis. Section~\ref{sec:safety} presents the
ICS-inspired intervention layer. Section~\ref{sec:experiments} reports
experiments across observable and degenerate geometries, and
Section~\ref{sec:conclusion} concludes.

%============================================================
% RELATED WORK
%============================================================
\section{Related Work}
\label{sec:related}

We organize prior work into four themes: (i) localization and position
holding in GNSS-denied environments, (ii) learned barrier functions defined
on raw perception, (iii) control synthesis directly in observation space, and
(iv) safe velocity filtering against nominal or adversarial commands. For each,
we clarify how our formulation differs, with particular care around the three
lines of work closest to ours.

\subsection{Position Holding under GNSS Denial}
The conventional route to station-keeping without GNSS is to run a
visual--inertial or LiDAR odometry front end and to close a position loop
around its estimate. These estimators are locally consistent but accumulate
unbounded drift in previously unvisited areas\cite{sprind2025}, and onboard
localization and mapping are prone to latency, noise, and drift that can lead
to estimation failure and, in turn, collision\cite{neuralcbf2024}. A large
body of work therefore focuses on improving the estimator---through sensor
fusion, cooperative localization, or terrain-relative
methods\cite{sprind2025}---so that a downstream position controller can hold
its setpoint reliably. Our work is orthogonal to this effort. Rather than
making the pose estimate more robust so that holding can depend on it, we
remove holding's dependence on pose altogether, providing a behavior that
remains valid precisely when the estimator is least trustworthy. In this sense
our method is complementary to, not competitive with, improved GNSS-denied
localization: it is intended as an independent, redundant layer that does not
share the estimator's failure modes.

\subsection{Learned Barrier Functions on Raw Perception}
A growing line of work synthesizes control barrier functions (CBFs) directly
from onboard sensing to avoid maintaining an explicit map. The neural CBF
safety filter of \cite{neuralcbf2024} operates in a local frame and does not
utilize a reconstructed map, relying instead on the instantaneous LiDAR scan.
This is close to our sensing philosophy, and we adopt a similar map-free
stance. The distinction is one of scope. In \cite{neuralcbf2024} the safety
guarantee is confined to the filter, while the underlying nominal controller
is driven by LiDAR-estimated odometry; the closed loop therefore remains
odometry-dependent, and yaw is regulated against an inertial frame. Our
pose-free property instead covers the entire translational governor, including
the behavior that plays the role of the nominal controller. Furthermore,
prior learned-CBF work targets collision avoidance around a functioning
navigation stack, whereas we study the more primitive problem of holding and
command response when that stack may be unavailable, and we contribute an
explicit account of when such measurement-space regulation is geometrically
feasible.

\subsection{Control Synthesis in Observation Space}
Closest to our sensing assumption is work that learns control certificates
directly in the robot's observation space. LOCUS learns barrier and Lyapunov
functions in observation space without assuming a separate perception-based
state estimator\cite{locus2023}, which aligns with our premise of regulating
on what the sensor measures rather than on a reconstructed state. We differ in
both problem and analysis. LOCUS is goal-driven and formulates navigation
toward a target, whereas we address the more primitive station-keeping regime
in which a null command is a first-class behavior rather than a degenerate
case of goal reaching. More importantly, observation-space synthesis has not,
to our knowledge, characterized \emph{when} such regulation is attainable: two
distinct positions may induce nearly identical observations, in which case no
policy can resist drift between them. Our range-Jacobian observability analysis
fills this gap by predicting the directions along which observation-space
holding must fail, turning a latent limitation of the observation-space
approach into an explicit, testable condition.

\subsection{Safe Velocity Filtering}
A final theme enforces safety by filtering a commanded velocity onto a safe
set. Composite CBF approaches combine multiple constraints and have been
validated against demanding inputs; notably, \cite{compositecbf2025}
demonstrates safety against a reference velocity issued by an adversarial
operator that intentionally attempts to collide the quadrotor. Reactive,
robot-centric filters similarly take a nominal velocity from a position
controller and output a filtered, safety-constrained velocity tracked by a
low-level controller\cite{reactive2026}. These methods are effective but are
position- or odometry-dependent and treat command filtering as the central
contribution. We deliberately do not position our command-intervention layer
as a novel filter. Our inevitable-collision-state-inspired criterion is
standard in spirit---it reasons about clearance, approach speed, and braking
feasibility---and its role in this paper is integrative: it operates on the
same raw range and inertial signals as the holding policy, so that both the
holding behavior and the command attenuation are produced without a
world-frame position estimate anywhere in the loop. To our knowledge, no prior
work provides a single command-conditioned, pose-free governor that unifies
measurement-space station-keeping, an observability account of its feasibility,
and local-sensing command intervention.

%============================================================
% COMPARISON TABLE + TRANSITION  (place at end of Sec. Related Work)
%============================================================
\subsection{Positioning Summary}
Table~\ref{tab:comparison} summarizes the preceding discussion along five axes
that jointly characterize a pose-free holding-and-response layer: whether the
method avoids an explicit map (\emph{map-free}); whether the actor that
produces motion---not merely a downstream filter---is free of pose and
odometry inputs (\emph{pose-free actor}); whether station-keeping under a null
command is treated as a first-class behavior (\emph{station-keeping}); whether
the feasibility of the approach is analyzed rather than assumed
(\emph{observability analysis}); and whether unsafe commands are rejected from
local sensing (\emph{command safety}). As the table shows, prior methods
satisfy subsets of these axes, but none combines a pose-free actor with an
explicit observability account of station-keeping---the gap this paper
addresses.

\begin{table}[t]
\centering
\caption{Positioning against representative prior work. A checkmark
(\checkmark) denotes full support, a tilde ($\sim$) partial support, and a
dash (--) no support. ``Pose-free actor'' requires the motion-generating
policy itself---not only a safety filter---to be free of pose/odometry inputs.}
\label{tab:comparison}
\renewcommand{\arraystretch}{1.3}
\setlength{\tabcolsep}{4pt}
\begin{tabular}{@{}lccccc@{}}
\toprule
\textbf{Method} &
\rotatebox{90}{Map-free} &
\rotatebox{90}{Pose-free actor} &
\rotatebox{90}{Station-keeping} &
\rotatebox{90}{Observability analysis} &
\rotatebox{90}{Command safety} \\
\midrule
GNSS-denied odometry + hold \cite{sprind2025} & -- & -- & \checkmark & -- & -- \\
Neural CBF filter \cite{neuralcbf2024}         & \checkmark & $\sim$ & -- & -- & \checkmark \\
LOCUS (obs.-space CBF/CLF) \cite{locus2023}    & \checkmark & \checkmark & -- & -- & \checkmark \\
Composite CBF \cite{compositecbf2025}          & -- & -- & -- & -- & \checkmark \\
Reactive robot-centric \cite{reactive2026}     & $\sim$ & -- & -- & -- & \checkmark \\
\midrule
\textbf{Ours}                                  & \checkmark & \checkmark & \checkmark & \checkmark & \checkmark \\
\bottomrule
\end{tabular}
\end{table}

\clearpage

%============================================================
% METHOD: Paper 1 velocity-command governor (final)
% Required packages: amsmath, amssymb, amsthm, bm, mathtools.
% Theorem envs required: definition, assumption, proposition, remark.
%============================================================

\section{Method}
\label{sec:formulation}
\label{sec:method}

We now describe the proposed measurement-space governor. The goal is not to
replace the entire flight stack. Instead, the learned module sits above a
conventional low-level velocity/attitude controller and produces a corrected
body-frame velocity command. The key restriction is imposed at the learned
module: the deployed governor is not given global pose, local odometry, a
reconstructed map, or an explicit velocity estimate. It must infer the
motion-relevant information needed for local holding and command response from
range--inertial history.

The resulting loop buffers raw LiDAR ranges and inertial cues over a short
horizon; a learned governor
maps this history and the operator command to a corrected velocity command, an
ICS-inspired intervention layer attenuates commands that are not locally
brakeable, and a conventional low-level controller tracks the final command.
We deliberately restrict the pose-free claim to the learned translational
governor. Attitude stabilization and low-level command execution are handled
by the flight controller and are not the object of this paper.

%============================================================
\subsection{System, Frames, and Information Constraint}
\label{sec:method_system}

Let the quadrotor state be
\begin{equation}
    x_t = (p_t, R_t, v_t, \omega_t),
    \label{eq:state}
\end{equation}
where $p_t \in \mathbb{R}^3$ is the position in an inertial world frame
$\mathcal{W}$, $R_t \in \mathrm{SO}(3)$ is the body-to-world rotation,
$v_t \in \mathbb{R}^3$ is the world-frame linear velocity, and
$\omega_t \in \mathbb{R}^3$ is the body angular velocity. This state is used
only for analysis, simulation reward construction, and evaluation. It is not
provided to the deployed learned governor.

We denote the body frame by $\mathcal{B}$. The operator velocity command is
expressed in a gravity-aligned body-yaw frame $\mathcal{Y}$ whose yaw follows
the vehicle heading while roll and pitch are removed. This convention avoids
requiring a global position frame and keeps commands intuitive for
teleoperation: forward, lateral, and vertical components correspond to the
current local vehicle heading. Let
\begin{equation}
    v^{\mathrm{cmd}}_t \in \mathcal{V} \subset \mathbb{R}^3
    \label{eq:cmd}
\end{equation}
be the operator command in $\mathcal{Y}$, with
$\|v^{\mathrm{cmd}}_t\| \leq v_{\max}$.

The learned governor receives only the following observation at time $t$:
\begin{equation}
    o_t = \bigl(r_t,\; y^{\mathrm{imu}}_t,\; v^{\mathrm{cmd}}_t,\;
    \tilde v_{t-1}\bigr),
    \label{eq:obs}
\end{equation}
where $r_t$ is the LiDAR range vector, $y^{\mathrm{imu}}_t$ denotes inertial
cues, and $\tilde v_{t-1}$ is the previously issued corrected velocity command.
The observation explicitly excludes
\begin{equation}
    p_t, \quad v_t, \quad \hat p_t, \quad \hat v_t, \quad
    \text{local odometry}, \quad \text{and reconstructed maps}.
    \label{eq:excluded_inputs}
\end{equation}
Because instantaneous translational velocity is not contained in a single
observation, the policy operates on a finite history window
\begin{equation}
    h_t = \{o_{t-H+1}, o_{t-H+2}, \ldots, o_t\},
    \label{eq:history}
\end{equation}
with window length $H$. This history can be realized by frame stacking, a
temporal convolutional encoder, or a recurrent encoder. In this paper we use a
fixed-window implementation as the default because it is stateless at
deployment and is less sensitive to hidden-state corruption under sensor
dropout or timing jitter.

%============================================================
\subsection{Range and Inertial Observation Model}
\label{sec:method_observation}

Consider a LiDAR with $N$ rays. The $i$-th ray has a known unit direction
$d_i \in \mathbb{S}^2$ expressed in the body frame, and its direction in the
world frame is $s_i = R_t d_i$. Let $\mathcal{O}$ denote the obstacle set. The
range return is
\begin{equation}
    r_{t,i} = \min\!\left(\rho(p_t, s_i, \mathcal{O}),\; r_{\max}\right)
    + \eta_{t,i},
    \qquad i=1,\ldots,N,
    \label{eq:range_model}
\end{equation}
where $\rho(p,s,\mathcal{O})$ is the distance from $p$ to the first intersection
with $\mathcal{O}$ along direction $s$, $r_{\max}$ is the maximum sensing range,
and $\eta_{t,i}$ models measurement noise. No-return rays are assigned the
saturated value $r_{\max}$ and are recorded by a validity mask
\begin{equation}
    m_{t,i} = \mathbb{1}[r_{t,i}<r_{\max}-\epsilon_r].
    \label{eq:valid_mask}
\end{equation}
We further associate to each beam a scalar reliability weight
$w_{t,i}\in[0,1]$, computed at deployment from return validity, range-noise
level, and neighbor consistency of the scan. This weight does \emph{not} require
surface normals and is used by the station-keeping loss and the
command-intervention layer to discount unreliable beams.

The inertial observation is
\begin{equation}
    y^{\mathrm{imu}}_t = \bigl(f^{\mathcal{B}}_t,\; \omega^{\mathcal{B}}_t,
    \; g^{\mathcal{B}}_t\bigr),
    \label{eq:imu_obs}
\end{equation}
where $f^{\mathcal{B}}_t$ is the specific force, $\omega^{\mathcal{B}}_t$ is the
body angular rate, and $g^{\mathcal{B}}_t$ is the gravity direction expressed in
$\mathcal{B}$. The gravity direction can be obtained from the attitude estimate
or a complementary filter; it provides roll/pitch stabilization information and
the rotation needed to express ray directions in the command frame
$\mathcal{Y}$, but it does not provide absolute position or translational
velocity.

%============================================================
\subsection{Measurement-Space Station-Keeping}
\label{sec:method_anchor}

Standard position holding requires a world-frame setpoint. We instead define
station-keeping directly in the sensor space. We use a hysteresis rule to enter
and exit the null-command (holding) regime, with an entry threshold $\epsilon_0$
and an exit threshold $\epsilon_1>\epsilon_0$:
\begin{equation}
    \chi_t =
    \begin{cases}
        1, & \chi_{t-1}=0 \ \text{and}\ \|v^{\mathrm{cmd}}_t\|\leq \epsilon_0,\\
        0, & \chi_{t-1}=1 \ \text{and}\ \|v^{\mathrm{cmd}}_t\|> \epsilon_1,\\
        \chi_{t-1}, & \text{otherwise},
    \end{cases}
    \label{eq:anchor_hysteresis}
\end{equation}
where $\chi_t\in\{0,1\}$ indicates that the holding anchor is active. On the
rising edge of $\chi_t$ we freeze the current range observation as an implicit
reference,
\begin{equation}
    r^\star = r_{t},\qquad m^\star = m_{t},
    \label{eq:range_anchor}
\end{equation}
and we retain $(r^\star,m^\star)$ until the anchor is cleared. The hysteresis
band $(\epsilon_0,\epsilon_1)$ prevents re-anchoring chatter near the command
threshold. The station-keeping objective is then to regulate subsequent
measurements toward this reference. We define the masked range-anchor error
\begin{equation}
    e^r_t = M_t \bigl(r_t-r^\star\bigr),
    \qquad
    M_t = \mathrm{diag}(m_t \odot m^\star \odot w_t),
    \label{eq:anchor_error}
\end{equation}
where $w_t=(w_{t,1},\ldots,w_{t,N})^\top$ are the beam reliability weights. In
practice we use a robust penalty
\begin{equation}
    \ell_{\mathrm{anch}}(t)
    = \sum_{i=1}^{N} M_{t,ii}\, \varphi_\delta(r_{t,i}-r^\star_i),
    \label{eq:anchor_loss}
\end{equation}
where $\varphi_\delta(\cdot)$ is the Huber loss. The robust form prevents a few
dynamic obstacles, dropouts, or range outliers from dominating the
station-keeping signal.

This definition is intentionally local. It does not claim that the vehicle
recovers an absolute position; it asks whether the vehicle can hold the same
exteroceptive view. If different translations produce nearly identical range
measurements, then no measurement-only controller can correct those
translations. This limitation is formalized next.

%============================================================
\subsection{Range-Jacobian Observability of Holding}
\label{sec:observability}

We analyze holding observability at a fixed attitude $R$; attitude motion is
delegated to the low-level controller and its first-order coupling into the
range map is neglected in this translational analysis. Let $g(p;R,\mathcal{O})$
denote the range measurement map,
\begin{equation}
    r = g(p;R,\mathcal{O}).
    \label{eq:measurement_map}
\end{equation}
For a small translation $\delta p$, the first-order range variation is
\begin{equation}
    \delta r \approx J(p)\,\delta p, \qquad
    J(p)=\frac{\partial g}{\partial p}(p;R,\mathcal{O}) \in \mathbb{R}^{N\times 3}.
    \label{eq:range_jacobian}
\end{equation}
For a ray that intersects a locally planar surface with outward normal $n_i$
and world-frame direction $s_i=R d_i$, the plane-intersection distance is
$r_i = n_i^\top(q_0-p)/(n_i^\top s_i)$ for a surface point $q_0$, and since the
denominator is independent of $p$,
\begin{equation}
    J_{i,:}(p)
    = -\frac{n_i^\top}{n_i^\top s_i},
    \qquad |n_i^\top s_i|>\epsilon_{\mathrm{grazing}}.
    \label{eq:range_jacobian_row}
\end{equation}
Thus, up to nonzero scalar factors induced by the incidence angle, the row
space of $J$ is spanned by the observed surface normals. Near-grazing rays are
excluded from this numerical analysis because the denominator in
\eqref{eq:range_jacobian_row} becomes ill-conditioned. We emphasize that surface
normals are used only for this offline observability analysis; the deployed
safety layer in Sec.~\ref{sec:safety} does not require them.

\begin{definition}[Translational measurement observability]
A unit translational direction $\xi \in \mathbb{R}^3$ is locally observable from
range measurements at $(p,R,\mathcal{O})$ if $\|J(p)\xi\|>0$. The configuration
is fully station-keepable in translation if $\mathrm{rank}\,J(p)=3$.
\end{definition}

\begin{proposition}[Aperture degeneracy]
If the normals of all valid range returns span a $k$-dimensional subspace of
$\mathbb{R}^3$ with $k<3$, then the unobservable translational subspace has
dimension at least $3-k$. Motions in this null space produce no first-order
change in range, and therefore cannot be corrected by any policy that only
observes range and inertial history.
\label{prop:aperture}
\end{proposition}

\begin{proof}
By \eqref{eq:range_jacobian_row}, every valid row of $J$ lies in the span of
the observed surface normals, up to nonzero scalar factors. Hence
$\mathrm{rank}\,J \leq k$. The rank--nullity theorem gives
$\dim\ker J = 3-\mathrm{rank}\,J \geq 3-k$. For any $\xi \in \ker J$, the
first-order measurement variation satisfies $\delta r \approx J\xi =0$, so range
measurements cannot distinguish motion along $\xi$ locally.
\end{proof}

This proposition explains common failure modes. Near a single flat wall,
translation tangent to the wall is weakly observable or unobservable. In a long
corridor, translation along the corridor axis produces little range change. In
open space, most rays return $r_{\max}$, $J\approx 0$, and station-keeping
degrades to inertial dead reckoning. We quantify this effect by reporting
station-keeping drift as a function of the conditioning of
\begin{equation}
    \mathcal{I}(p) = J(p)^\top W_t J(p),
    \qquad
    \sigma_{\min}(p) = \sqrt{\lambda_{\min}(\mathcal{I}(p))},
    \label{eq:information_matrix}
\end{equation}
where $W_t=\mathrm{diag}(m_t\odot w_t)$ contains the same validity and
reliability weights used in \eqref{eq:anchor_error}. In simulation, $J$ can be
computed from known hit normals or by finite differences of the ray-casting
function. In real-robot experiments, where surface normals are not always
reliable, we report a proxy conditioning score from local scan geometry and
interpret it as an empirical observability indicator rather than as an exact
Jacobian.

%============================================================
\subsection{Range--Inertial History Governor}
\label{sec:governor}

The learned governor maps the history $h_t$ and the active anchor state to a
corrected velocity command. Its output is not an unconstrained velocity vector.
We parameterize the command as an attenuation of the operator input plus a
bounded corrective component:
\begin{equation}
    v^{g}_t
    = \alpha_t v^{\mathrm{cmd}}_t + v^{\mathrm{corr}}_t,
    \qquad
    \alpha_t \in [0,1], \quad
    \|v^{\mathrm{corr}}_t\| \leq v_{\mathrm{corr,max}}.
    \label{eq:governor_output}
\end{equation}
The scalar $\alpha_t$ represents how much of the operator command is preserved:
close to one when the scene is safe and observable, and smaller when the command
is unsafe or incompatible with holding. The corrective component
$v^{\mathrm{corr}}_t$ lets the policy counteract local drift under a null command
and bias the vehicle away from local hazards under a nonzero command.

A typical implementation is
\begin{align}
    z^r_t &= E_r(r_{t-H+1:t}, m_{t-H+1:t}, w_{t-H+1:t}), \\
    z^i_t &= E_i(y^{\mathrm{imu}}_{t-H+1:t}), \\
    z^c_t &= E_c(v^{\mathrm{cmd}}_{t-H+1:t}, \tilde v_{t-H:t-1}), \\
    z_t   &= \mathrm{Fuse}(z^r_t,z^i_t,z^c_t,e^r_t,\chi_t), \\
    (\alpha_t,v^{\mathrm{corr}}_t) &= \Pi_\theta(z_t),
    \label{eq:network_architecture}
\end{align}
where $E_r$ is a range encoder, $E_i$ is an inertial encoder, $E_c$ is a command
history encoder, $\chi_t$ is the anchor flag from \eqref{eq:anchor_hysteresis},
and $\Pi_\theta$ is an MLP policy head. The range encoder can be a small
convolution over beam angle and time; the temporal encoder can be a
stacked-window MLP or a temporal convolutional network. The governor command is
clipped to the admissible velocity set:
\begin{equation}
    \bar v^g_t = \mathrm{clip}_{\mathcal{V}}(v^{g}_t).
    \label{eq:clip_governor}
\end{equation}
This command is subsequently passed through the braking-feasibility intervention
layer of Sec.~\ref{sec:safety}.

\begin{remark}[Actor-level information restriction]
The policy may be trained with privileged simulation rewards and an asymmetric
critic, but its deployed input is exactly $h_t$ and the anchor state. No pose,
odometry, explicit velocity estimate, or map feature is passed to $\Pi_\theta$.
\end{remark}

%============================================================
\subsection{ICS-Inspired Command Intervention}
\label{sec:safety}

The learned governor should not obey a nonzero command if the resulting motion
is not locally brakeable. We therefore add a deterministic intervention layer
inspired by inevitable collision states (ICS). It is not an exact reachable-set
computation but a conservative command-level surrogate based on a
\emph{line-of-sight point-obstacle model}, observed range dynamics, and the
maximum deceleration the low-level controller can realize. Importantly, the
layer depends only on line-of-sight range and ray direction; it does not use
surface normals, so it does not inherit the normal-estimation fragility of real
LiDAR scans.

\paragraph{Line-of-sight point model.}
Each valid return defines a momentary point obstacle
$q_{t,i} = p_t + r_{t,i}\,s_i$ at line-of-sight distance $r_{t,i}$. For a
candidate velocity $u\in\mathcal{V}$, the rate at which the sensor approaches
this fixed point is, exactly to first order, the line-of-sight component
\begin{equation}
    c_i(u) = \bigl[\,d_i^{\mathcal{Y}\top} u\,\bigr]_+ ,
    \label{eq:commanded_closing_ray}
\end{equation}
where $d_i^{\mathcal{Y}}$ is the ray direction expressed in the command frame
$\mathcal{Y}$ (obtained from $g^{\mathcal{B}}_t$ and the vehicle yaw). Unlike a
radial \emph{approximation}, $c_i(u)$ is the \emph{exact} approach rate to the
currently observed point at distance $r_{t,i}$, which is the quantity relevant
to per-point collision. We evaluate $c_i$ only on valid, reliable rays
($m_{t,i}=1$, $w_{t,i}>w_{\min}$).

\paragraph{Observed closing indicator.}
For each valid ray we estimate a range-rate from the history window by a
recency-weighted least-squares line fit, jointly optimizing slope and intercept,
\begin{equation}
    (\hat a_{t,i},\,\hat b_{t,i})
    = \arg\min_{a,b\in\mathbb{R}}
    \sum_{k=0}^{H_r-1} \omega_k
    \bigl(r_{t-k,i} - b - a(-k\Delta t)\bigr)^2,
    \qquad
    \dot r^{H}_{t,i} = \hat a_{t,i},
    \label{eq:range_rate_fit}
\end{equation}
with recency weights $\omega_k$, and define the per-ray observed closing
indicator
\begin{equation}
    \hat v^{\mathrm{obs}}_{i,t} = \bigl[-\dot r^{H}_{t,i}\bigr]_+ .
    \label{eq:observed_closing_speed}
\end{equation}
For oblique surfaces the range-rate can exceed the true point-approach rate
because the hit point slides along the surface; we therefore use
$\hat v^{\mathrm{obs}}_{i,t}$ only as a \emph{conservative} closing indicator
(for the latency margin and risk logging), never to relax the intervention. It
requires no external velocity estimate.

\paragraph{Braking-feasibility attenuation.}
Given a deceleration bound $a_{\max}$ and safety margin $d_{\mathrm{safe}}$, the
braking-distance speed limit for ray $i$ is
\begin{equation}
    v_{\mathrm{safe}}(r)
    = \sqrt{2a_{\max}\,[r-d_{\mathrm{safe}}]_+}.
    \label{eq:vsafe}
\end{equation}
To account for processing and actuation delay, we shrink the available clearance
per ray by
\begin{equation}
    r^{\tau}_{t,i}
    = r_{t,i} - \tau_{\mathrm{lat}}\,
      \max\{\hat v^{\mathrm{obs}}_{i,t},\; c_i(\bar v^g_t)\},
    \label{eq:latency_margin}
\end{equation}
where $\tau_{\mathrm{lat}}$ is a conservative latency bound. The command
attenuation factor is
\begin{equation}
    \beta_t
    = \min\!\left(1,
    \min_{\substack{i:\,m_{t,i}=1,\ w_{t,i}>w_{\min}\\
    c_i(\bar v^g_t)>\epsilon_c}}
    \frac{v_{\mathrm{safe}}(r^{\tau}_{t,i})}
         {c_i(\bar v^g_t)+\epsilon_c}
    \right),
    \label{eq:beta_attenuation}
\end{equation}
with a small $\epsilon_c>0$. The final velocity command issued to the low-level
controller is
\begin{equation}
    \tilde v_t = \beta_t \bar v^g_t + (1-\beta_t)\,v^{\mathrm{brake}}_t,
    \label{eq:final_command}
\end{equation}
where $v^{\mathrm{brake}}_t$ is a conservative fallback command satisfying
$d_i^{\mathcal{Y}\top} v^{\mathrm{brake}}_t \le 0$ on every valid reliable ray,
i.e., the fallback never increases the line-of-sight approach to any observed
point. In the simplest case $v^{\mathrm{brake}}_t=0$; in cluttered scenes it may
include a small reverse or vertical component selected from the most open
sector, subject to the same non-approach condition. If
$d_t=\min_i r_{t,i}<d_{\mathrm{emg}}$, the system bypasses the learned output and
issues an emergency stop or pilot takeover.

The risk indicator used for logging and training is
\begin{equation}
    \rho_t
    = \max_{i:\,m_{t,i}=1,\ w_{t,i}>w_{\min}}
    \bigl[\,\hat v^{\mathrm{obs}}_{i,t} - v_{\mathrm{safe}}(r_{t,i})\bigr]_+ ,
    \label{eq:risk_indicator}
\end{equation}
which becomes positive when an observed closing indicator exceeds the local
braking-feasibility bound. This is an ICS-inspired surrogate: it flags that the
vehicle may be entering a region in which collision becomes unavoidable unless
the command is attenuated.

\begin{assumption}[Velocity-command execution]
The low-level controller can reduce the component of vehicle velocity toward an
observed obstacle with deceleration at least $a_{\max}$ after a delay bounded by
$\tau_{\mathrm{lat}}$, within the command and attitude limits used in the
experiments.
\label{assumption:execution}
\end{assumption}

\begin{assumption}[Observation sufficiency and quasi-static obstacles]
Over the braking horizon, collision-relevant obstacles are quasi-static, and
every such obstacle produces at least one valid, reliable return
($m_{t,i}=1$, $w_{t,i}>w_{\min}$), so that no collision-relevant surface lies
entirely between adjacent beams or in an occluded region.
\label{assumption:observation}
\end{assumption}

\begin{proposition}[Command-level braking feasibility]
Under Assumptions~\ref{assumption:execution} and \ref{assumption:observation},
and with $v^{\mathrm{brake}}_t$ satisfying the non-approach condition above, the
final command \eqref{eq:final_command} obeys
$c_i(\tilde v_t) \leq v_{\mathrm{safe}}(r^{\tau}_{t,i})$ for every valid,
reliable ray. Consequently, the commanded motion does not ask the vehicle to
approach any observed point obstacle faster than the speed at which the
low-level controller can brake within the available clearance.
\label{prop:command_feasibility}
\end{proposition}

\begin{proof}
Fix a valid reliable ray $i$. By convexity and nonnegativity of $[\cdot]_+$ and
$\beta_t\in[0,1]$,
\begin{align*}
    c_i(\tilde v_t)
    &= \bigl[d_i^{\mathcal{Y}\top}(\beta_t \bar v^g_t
       + (1-\beta_t)v^{\mathrm{brake}}_t)\bigr]_+ \\
    &\leq \beta_t\bigl[d_i^{\mathcal{Y}\top}\bar v^g_t\bigr]_+
       + (1-\beta_t)\bigl[d_i^{\mathcal{Y}\top}v^{\mathrm{brake}}_t\bigr]_+ \\
    &= \beta_t\, c_i(\bar v^g_t),
\end{align*}
where the last equality uses the non-approach condition
$d_i^{\mathcal{Y}\top}v^{\mathrm{brake}}_t\le 0$. If
$c_i(\bar v^g_t)\le\epsilon_c$ the claim is immediate; otherwise, by
\eqref{eq:beta_attenuation},
$\beta_t \le v_{\mathrm{safe}}(r^{\tau}_{t,i})/(c_i(\bar v^g_t)+\epsilon_c)$,
hence $\beta_t c_i(\bar v^g_t) \le v_{\mathrm{safe}}(r^{\tau}_{t,i})$. By
\eqref{eq:vsafe}, $v_{\mathrm{safe}}(r^{\tau}_{t,i})$ is the largest closing
speed brakeable within clearance $[r^{\tau}_{t,i}-d_{\mathrm{safe}}]_+$ under
deceleration $a_{\max}$; Assumption~\ref{assumption:execution} makes this
deceleration realizable, and Assumption~\ref{assumption:observation} ensures the
checked rays cover all collision-relevant obstacles.
\end{proof}

\begin{remark}[Difference from CBF filtering]
The intervention layer is not a new control barrier function. It does not solve
a state-based QP and does not require a pose or velocity estimate as policy
input. It is a command-level braking-feasibility check coupled to the
measurement-space governor. CBF-style filters serve as related work and
baselines, but the contribution here is the pose-free governor and its
measurement-space holding behavior.
\end{remark}

%============================================================
\subsection{Training Objective}
\label{sec:training_objective}

The policy is trained in simulation with privileged quantities used only for
reward construction and critic learning; these are not part of the deployed
governor input. The objective combines command tracking, measurement-space
holding, smoothness, safety, and a light penalty on relying on the intervention
layer:
\begin{equation}
    R_t
    = w_{\mathrm{trk}} R^{\mathrm{trk}}_t
    + w_{\mathrm{anch}} R^{\mathrm{anch}}_t
    + w_{\mathrm{safe}} R^{\mathrm{safe}}_t
    + w_{\mathrm{int}} R^{\mathrm{int}}_t
    - w_{\mathrm{cmp}} R^{\mathrm{cmp}}_t
    - w_{\mathrm{sm}} \|\tilde v_t-\tilde v_{t-1}\|^2.
    \label{eq:reward_total}
\end{equation}
The tracking term is active under nonzero commands:
\begin{equation}
    R^{\mathrm{trk}}_t
    = \mathbb{1}[\chi_t=0]\,
    \exp\!\left(-\lambda_v \|v^{\mathcal{Y}}_t - v^{\mathrm{cmd}}_t\|^2\right),
    \label{eq:tracking_reward}
\end{equation}
where $v^{\mathcal{Y}}_t$ is the true velocity in the command frame, available
in simulation. The anchor term is active under a null command:
\begin{equation}
    R^{\mathrm{anch}}_t
    = -\mathbb{1}[\chi_t=1]\,\ell_{\mathrm{anch}}(t).
    \label{eq:anchor_reward}
\end{equation}
The command-compliance penalty is weighted by the braking-feasibility factor
$\beta_t$ so that it never fights the intervention layer: the more the governor
command is unsafe (small $\beta_t$), the less the policy is pushed to reproduce
the operator command,
\begin{equation}
    R^{\mathrm{cmp}}_t
    = \beta_t\,\|\bar v^g_t-v^{\mathrm{cmd}}_t\|^2 .
    \label{eq:compliance_reward}
\end{equation}
Because $\beta_t$ is derived from the commanded closing speed rather than the
observed one, this gate reflects the danger of the command itself and remains
consistent with \eqref{eq:beta_attenuation}. The safety term penalizes margin
violation and the ICS-inspired risk surrogate:
\begin{equation}
    R^{\mathrm{safe}}_t
    = -\mathbb{1}[d_t<d_{\mathrm{safe}}]
      -\kappa_{\rho}\rho_t
      -C_{\mathrm{coll}}\mathbb{1}[\mathrm{collision}].
    \label{eq:safety_reward}
\end{equation}
Finally, the intervention-usage penalty discourages the policy from emitting
commands that must be repeatedly rescued by the deterministic layer:
\begin{equation}
    R^{\mathrm{int}}_t = -(1-\beta_t)^2 .
    \label{eq:intervention_reward}
\end{equation}

\begin{remark}[Role of the intervention-usage penalty]
The penalty \eqref{eq:intervention_reward} is a soft shaping term, not a safety
mechanism: it pushes the required attenuation into the learned gain $\alpha_t$ so
that $\bar v^g_t$ is already brake-feasible ($\beta_t\to1$), turning
$\Pi_\theta$ into a genuine safety governor rather than an ``emit-and-filter''
policy. The hard guarantee of
Proposition~\ref{prop:command_feasibility} is provided by the deterministic layer
regardless of this term. Its weight $w_{\mathrm{int}}$ is kept modest so that, in
genuinely constrained situations (e.g., an operator commanding straight into a
wall), the correct persistent attenuation $\beta_t<1$ is not over-penalized.
\end{remark}

We use a reward-penalty formulation for the main system rather than a Lagrangian
constraint because the purpose of the paper is the measurement-space holding
formulation and its observability limits, not a new constrained-RL algorithm. A
Lagrangian or PID-Lagrangian variant can be reported as an ablation.

The critic may use privileged simulator state
\begin{equation}
    s^{\mathrm{priv}}_t = (p_t,R_t,v_t,\omega_t,d_t,\rho_t,\beta_t,
    \mathrm{collision})
    \label{eq:privileged_state}
\end{equation}
for value estimation during training. At deployment, only the actor
$\Pi_\theta$ and the deterministic intervention layer are retained.

Optionally, we add an auxiliary head from the shared feature $z_t$ to estimate
body-frame velocity or range-rate,
\begin{equation}
    \hat v^{\mathrm{aux}}_t = A_\psi(z_t),
    \qquad
    L_{\mathrm{aux}} = \|\hat v^{\mathrm{aux}}_t - v^{\mathcal{Y}}_t\|^2,
    \label{eq:auxiliary_loss}
\end{equation}
using simulator ground truth. This loss shapes the history representation but
its output is not fed back as a policy input; the auxiliary head can be removed
at deployment and is evaluated as an ablation, not assumed by the safety
mechanism.

%============================================================
\subsection{Deployment Loop}
\label{sec:deployment_loop}

At runtime, the module performs the following at the policy rate:
\begin{enumerate}
    \item read $r_t$, mask $m_t$, reliability weights $w_t$, IMU cues
    $y^{\mathrm{imu}}_t$, and command $v^{\mathrm{cmd}}_t$;
    \item update the fixed history buffer $h_t$;
    \item update the anchor flag $\chi_t$ by the hysteresis rule
    \eqref{eq:anchor_hysteresis}; on a rising edge, set $r^\star=r_t$,
    $m^\star=m_t$;
    \item compute the governor command $\bar v^g_t$ via
    \eqref{eq:network_architecture}--\eqref{eq:clip_governor};
    \item compute the attenuation $\beta_t$ from local range history via
    \eqref{eq:range_rate_fit}--\eqref{eq:beta_attenuation};
    \item issue the final command $\tilde v_t$ in \eqref{eq:final_command} to
    the low-level velocity/attitude controller.
\end{enumerate}
The module logs $d_t$, $\rho_t$, $\beta_t$, $\ell_{\mathrm{anch}}$, and, in
simulation, $\sigma_{\min}$ from \eqref{eq:information_matrix}. These signals
separate policy failure from geometric unobservability: high anchor error under
well-conditioned geometry indicates a policy or controller failure, whereas
drift aligned with $\ker J$ under poor conditioning is an expected limitation of
measurement-only holding.

%============================================================
\subsection{Guarantees and Limitations}
\label{sec:guarantees}

The method provides conditional, not absolute, guarantees.
Measurement-space station-keeping can only regulate directions visible through
the range map: if $\xi\in\ker J(p)$, range measurements carry no first-order
information about motion along $\xi$, and any governor without additional
position or velocity feedback must rely on inertial cues that drift. The
observability analysis is first-order and fixed-attitude; large attitude
excursions during holding are outside its scope. Command intervention is
meaningful only for obstacles inside the sensing range and under
Assumptions~\ref{assumption:execution}--\ref{assumption:observation}; in
particular, thin obstacles that fall between adjacent beams or occluded surfaces
are not covered by Proposition~\ref{prop:command_feasibility}, and the observed
range-rate indicator can over-flag on oblique surfaces. These limitations are
not treated as failures of the method; they define the physically attainable
regime and are explicitly tested in the experiments.

\clearpage

%============================================================
% EXPERIMENTAL PROTOCOL AND EVALUATION MATRIX
% Paper 1: velocity-command governor version
% Designed to be included as an IEEE-style experiment section.
% Required packages in the main file: amsmath, amssymb, booktabs, array.
% Optional but recommended: multirow, makecell, tabularx.
%============================================================

%============================================================
\section{Experimental Protocol}
\label{sec:experiments}
%============================================================

The experiments are designed to test the four claims of the proposed
measurement-space governor: (i) the learned governor operates without pose,
odometry, map, or explicit velocity inputs; (ii) station-keeping can be posed as
regulation in LiDAR measurement space; (iii) the success and failure of such
regulation are predicted by range observability; and (iv) nonzero operator
commands can be attenuated when they violate local braking feasibility.  The
same experimental protocol is used in simulation and, where feasible, on the
real quadrotor.  External motion capture or localization is allowed only for
logging and evaluation; it is never supplied to the deployed governor.

For compactness, we refer to the paper's main contributions as follows.
\begin{itemize}
    \item \textbf{C1:} pose-, odometry-, map-, and explicit-velocity-free
    learned-governor formulation;
    \item \textbf{C2:} measurement-space station-keeping under a null command;
    \item \textbf{C3:} range-Jacobian observability analysis and predictable
    drift in degenerate geometries;
    \item \textbf{C4:} ICS-inspired command intervention under nonzero
    commands.
\end{itemize}

%============================================================
\subsection{Scenarios}
\label{sec:exp_scenarios}
%============================================================

We evaluate the governor in scenarios chosen to separate controller performance
from measurement observability.  The first group contains controlled geometric
scenes for station-keeping.  The second group contains command-tracking scenes
with increasing safety pressure.  The third group transfers representative
cases to the physical platform.

\begin{table*}[t]
\centering
\caption{Evaluation scenarios.  The geometric scenarios are chosen to span the
conditioning of the range Jacobian, from well-conditioned corners to open-field
no-return observations.}
\label{tab:scenario_matrix}
\small
\setlength{\tabcolsep}{3pt}
\begin{tabular}{L{0.08\textwidth}L{0.18\textwidth}L{0.28\textwidth}L{0.30\textwidth}L{0.08\textwidth}}
\toprule
\textbf{ID} & \textbf{Scenario} & \textbf{Purpose} & \textbf{Expected property} & \textbf{C} \\
\midrule
S1 & Corner / room corner & Test station-keeping in an observable local geometry. & Surface normals span multiple directions; range Jacobian is well conditioned for local translation. & C2, C3 \\
S2 & Single flat wall & Test aperture degeneracy. & Motion normal to the wall is observable; motion tangent to the wall is weakly observable or unobservable. & C3 \\
S3 & Corridor / tunnel & Test drift along a dominant axis. & Lateral motion is observable, but translation along the corridor axis changes ranges weakly. & C3 \\
S4 & Open field / no-return space & Test the failure boundary. & Most beams return the maximum range; the range Jacobian is close to zero and station-keeping reduces to inertial dead reckoning. & C3 \\
S5 & Cluttered indoor scene & Test normal command tracking and local obstacle interaction. & Multiple surfaces are visible; local range variation should provide sufficient motion cues. & C1, C2, C4 \\
S6 & Wall-approach command & Test unsafe command attenuation. & The operator command points toward the nearest obstacle; braking feasibility should trigger attenuation before impact. & C4 \\
S7 & Gap / doorway passage & Test non-overconservative behavior. & The governor should preserve command tracking when clearance and braking authority are sufficient. & C4 \\
S8 & Dynamic obstacle crossing & Test robustness to changing local clearance. & The closest obstacle and approach direction change over time; intervention should be smooth and reactive. & C4 \\
S9 & Real indoor corner / corridor & Validate station-keeping and observability trends on hardware. & Real LiDAR noise, latency, airflow, and controller dynamics are present. & C1--C3 \\
S10 & Real unsafe-command approach & Validate command attenuation on hardware. & Operator commands the vehicle toward an obstacle; safety margin and manual override are enforced. & C4 \\
\bottomrule
\end{tabular}
\end{table*}

The controlled geometric scenes S1--S4 are not intended to make station-keeping
easier.  They are designed to expose when the information available to a
range-only governor is sufficient or insufficient.  In particular, poor
performance in S4 is not counted as a method failure if it is consistent with
low range observability; it is evidence for the claimed observability boundary.

%============================================================
\subsection{Baselines and Ablations}
\label{sec:exp_baselines}
%============================================================

The baselines are selected to isolate the effect of each component rather than
to maximize the number of comparisons.  All learned variants use the same
low-level velocity/attitude controller.  The difference lies only in the inputs
available to the learned governor, the use of the measurement-space anchor, the
history encoder, and the intervention rule.

\begin{table*}[t]
\centering
\caption{Baselines and ablations.  Variants B0--B4 test the information
constraint and station-keeping mechanism; B5--B8 test command intervention and
safety behavior.}
\label{tab:baseline_matrix}
\small
\setlength{\tabcolsep}{3pt}
\begin{tabular}{L{0.08\textwidth}L{0.22\textwidth}L{0.37\textwidth}L{0.22\textwidth}}
\toprule
\textbf{ID} & \textbf{Method} & \textbf{Description} & \textbf{Question answered} \\
\midrule
B0 & Direct velocity execution & The operator command is sent directly to the low-level controller: $\tilde v_t=v^{\mathrm{cmd}}_t$.  No learned governor, anchor, or intervention is used. & What does the underlying velocity controller do by itself? \\
B1 & State/velocity-aided governor & Same architecture as ours, but explicit velocity and/or local odometry are provided to the governor.  This is an upper-bound reference, not a deployable variant under our information constraint. & How much is lost by removing explicit state inputs? \\
B2 & Ours without history & The governor receives only the current observation $o_t$ rather than the window $h_t$. & Is temporal information necessary for motion inference? \\
B3 & Ours without range anchor & The learned governor is trained for command tracking but the station-keeping anchor loss is removed under a null command. & Does measurement-space anchoring reduce drift? \\
B4 & Ours without IMU cues & The governor uses range history and commands but removes inertial cues. & Do inertial cues improve short-horizon motion inference and robustness? \\
B5 & Ours without intervention & The learned governor is used, but the ICS-inspired attenuation layer is disabled. & Does the safety layer reduce unsafe-command failures? \\
B6 & Distance-only filter & A deterministic filter attenuates commands only when $d_t<d_{\mathrm{warn}}$, without using approach speed or braking feasibility. & Is distance-only reactivity sufficient? \\
B7 & Braking-distance heuristic & A non-learned command governor applies the same braking-feasibility attenuation using range-rate estimates, but without the learned correction term. & What is gained by learning beyond a hand-designed rule? \\
B8 & Simplified CBF/VO-style filter & A model-based safety-filter baseline, if implemented, modifies the command according to a local distance or velocity-obstacle constraint. & How does the proposed governor compare with a standard safety-filter style baseline? \\
Ours & Proposed governor & Range--inertial history, measurement-space anchor, learned command correction, and ICS-inspired intervention. & Full method. \\
\bottomrule
\end{tabular}
\end{table*}

B1 is used only as a diagnostic upper bound.  It is not treated as a fair
competitor under the paper's information constraint, because it receives the
explicit velocity or odometry signal that our deployed governor is designed to
avoid.  B8 is optional because a faithful reproduction of recent neural or
composite CBF systems would require different state assumptions and obstacle
representations; the baseline is used only to show how a conventional local
filter behaves under the same test scenes.

%============================================================
\subsection{Metrics}
\label{sec:exp_metrics}
%============================================================

The metrics are grouped into station-keeping, observability, tracking, safety,
and deployment categories.  Quantities involving ground-truth position or
velocity are used only for evaluation and training diagnostics, not as policy
inputs.

\subsubsection{Station-keeping metrics}
For a null-command interval beginning at $t_0$ and ending at $T$, the physical
drift is measured by
\begin{equation}
    D_{\mathrm{pos}}(T) = \|p_T-p_{t_0}\|_2,
    \label{eq:metric_position_drift}
\end{equation}
where $p_t$ is obtained from simulation ground truth or motion capture.  The
measurement-space drift is
\begin{equation}
    D_{r}(T) =
    \sqrt{\frac{1}{|\mathcal{T}|}\sum_{t\in\mathcal{T}}
    \frac{\|M_t\odot(r_t-r^\star)\|_2^2}{\|M_t\|_1+\epsilon}},
    \label{eq:metric_range_drift}
\end{equation}
where $M_t$ is the valid-hit mask and $r^\star$ is the anchor range pattern.
The projection of the drift onto weakly observable directions is also reported
when the right singular vectors of the range Jacobian are available:
\begin{equation}
    D_{\mathrm{null}}(T) = \|V_{\mathrm{weak}}^\top(p_T-p_{t_0})\|_2.
    \label{eq:metric_null_drift}
\end{equation}

\subsubsection{Observability metrics}
Let $J_t$ be the local range Jacobian computed from hit rays or from the known
scene geometry in simulation.  We report
\begin{equation}
    \sigma_{\min}(J_t), \qquad
    \kappa(J_t)=\frac{\sigma_{\max}(J_t)}{\sigma_{\min}(J_t)+\epsilon},
    \label{eq:metric_observability}
\end{equation}
with rank estimated by thresholding singular values.  The main validation is a
correlation between station-keeping drift and the conditioning of $J_t$ across
scenes and initial poses.  We report both binned averages and a rank or
Spearman correlation between $D_{\mathrm{pos}}$ and
$1/(\overline{\sigma}_{\min}+\epsilon)$.

\subsubsection{Command-tracking metrics}
For nonzero commands, the tracking error is evaluated as
\begin{equation}
    E_{\mathrm{trk}} =
    \sqrt{\frac{1}{T}\sum_{t=1}^{T}
    \|v^{\mathcal{Y}}_t-v^{\mathrm{cmd}}_t\|_2^2},
    \label{eq:metric_tracking_rmse}
\end{equation}
where $v^{\mathcal{Y}}_t$ is the measured or simulated vehicle velocity in the
body-yaw frame.  We also report the error with respect to the corrected command
$\tilde v_t$ to separate command attenuation from low-level tracking error.
Command preservation is measured by
\begin{equation}
    P_{\mathrm{cmd}} =
    \frac{1}{T}\sum_{t=1}^{T}
    \frac{\tilde v_t^\top v^{\mathrm{cmd}}_t}
    {\|\tilde v_t\|_2\|v^{\mathrm{cmd}}_t\|_2+\epsilon},
    \label{eq:metric_command_preservation}
\end{equation}
and smoothness by
\begin{equation}
    S_{\mathrm{cmd}} = \frac{1}{T-1}\sum_{t=2}^{T}
    \|\tilde v_t-\tilde v_{t-1}\|_2^2.
    \label{eq:metric_smoothness}
\end{equation}

\subsubsection{Safety and intervention metrics}
The minimum clearance is
\begin{equation}
    d_{\min}^{\mathrm{episode}} = \min_t d_t,
    \label{eq:metric_min_clearance}
\end{equation}
and the collision rate is the fraction of trials in which the vehicle intersects
an obstacle or violates a hardware safety boundary.  The ICS-inspired violation
ratio is
\begin{equation}
    R_{\mathrm{ICS}} =
    \frac{1}{T}\sum_{t=1}^{T}\mathbf{1}
    \left[v_{\parallel,t}^{\mathrm{range}}>v_{\mathrm{safe}}(d_t)\right],
    \label{eq:metric_ics_ratio}
\end{equation}
where $v_{\parallel,t}^{\mathrm{range}}$ is estimated from range history.  The
intervention rate is
\begin{equation}
    R_{\mathrm{int}} =
    \frac{1}{T}\sum_{t=1}^{T}\mathbf{1}
    \left[\|\tilde v_t-v^{\mathrm{cmd}}_t\|_2>\epsilon_{\mathrm{int}}\right],
    \label{eq:metric_intervention_rate}
\end{equation}
and the attenuation magnitude is
\begin{equation}
    A_{\mathrm{cmd}} =
    \frac{1}{T}\sum_{t=1}^{T}
    \frac{\|v^{\mathrm{cmd}}_t\|_2-\|\tilde v_t\|_2}
    {\|v^{\mathrm{cmd}}_t\|_2+\epsilon}.
    \label{eq:metric_attenuation}
\end{equation}

\subsubsection{Information-constraint audit}
To support C1, each deployed policy is audited by logging its input tensor and
verifying that no pose, odometry, explicit velocity, reconstructed map, or
motion-capture signal is included.  We report this audit as a binary checklist
for each method variant.  Baselines B1 and state-aided variants are explicitly
marked as using privileged inputs.

%============================================================
\subsection{Experiment Matrix}
\label{sec:exp_matrix}
%============================================================

Table~\ref{tab:experiment_matrix} maps each experiment to scenarios, baselines,
metrics, expected outcomes, and the corresponding contributions.  The expected
outcomes are stated as qualitative hypotheses; final numerical thresholds are
set after simulator calibration and hardware safety tests.

\begin{table*}[p]
\centering
\caption{Experiment matrix linking scenarios, baselines, metrics, expected
outcomes, and contributions.}
\label{tab:experiment_matrix}
\scriptsize
\renewcommand{\arraystretch}{0.82}
\setlength{\tabcolsep}{2pt}
\begin{tabular}{L{0.07\textwidth}L{0.15\textwidth}L{0.15\textwidth}L{0.20\textwidth}L{0.28\textwidth}L{0.06\textwidth}}
\toprule
\textbf{Exp.} & \textbf{Scenario(s)} & \textbf{Baselines} & \textbf{Metrics} & \textbf{Expected result} & \textbf{C} \\
\midrule
E1 & S1--S4: corner, wall, corridor, open field under $v^{\mathrm{cmd}}=0$ & B0, B2, B3, B4, Ours & $D_{\mathrm{pos}}$, $D_r$, $D_{\mathrm{null}}$, episode stability & Ours reduces drift in observable scenes relative to B0/B3/B4; it does not artificially hide failure in open field.  Explicit-state B1/B2 provides an upper bound. & C1--C3 \\
E2 & S1--S4 with randomized pose, wall distance, corridor width, sensor range, and valid-beam density & Ours, B3, B4 & $\sigma_{\min}(J)$, rank estimate, $D_{\mathrm{pos}}$, drift direction, drift--observability correlation & Station-keeping drift increases as the range Jacobian becomes ill conditioned; weak-direction drift aligns with the approximate null space of $J$. & C3 \\
E3 & S5 with step, sinusoidal, and piecewise-constant nonzero commands in observable geometry & B0, B1, B2, B4, Ours & $E_{\mathrm{trk}}$, $P_{\mathrm{cmd}}$, $S_{\mathrm{cmd}}$, intervention rate & Ours tracks feasible commands with moderate degradation relative to state-aided variants, while outperforming no-history variants. & C1, C2 \\
E4 & S6: commanded motion toward a wall or obstacle at increasing speed & B0, B5, B6, B7, Ours & collision rate, $d_{\min}^{\mathrm{episode}}$, $R_{\mathrm{ICS}}$, $R_{\mathrm{int}}$, $A_{\mathrm{cmd}}$ & Ours attenuates before the braking-feasibility boundary is violated, reducing collisions relative to B0/B5 and intervening earlier than distance-only filtering. & C4 \\
E5 & S7: command through a gap / doorway with sufficient clearance & B5, B6, B7, Ours & $E_{\mathrm{trk}}$, $P_{\mathrm{cmd}}$, $R_{\mathrm{int}}$, minimum clearance & Ours preserves feasible commands and is less conservative than distance-only filtering while maintaining clearance. & C4 \\
E6 & S8: moving obstacle crossing and changing nearest-obstacle identity & B5, B6, B7, Ours & collision rate, $R_{\mathrm{ICS}}$, command smoothness, minimum clearance & Ours reacts smoothly to changing local risk; intervention does not produce oscillatory command switching. & C4 \\
E7 & Information ablation across S1--S6 & B1, B2, B3, B4, Ours & performance gap to state-aided upper bound, input-audit checklist, tracking and drift metrics & Explicit velocity/odometry improves the upper bound, but the proposed governor remains functional under the strict actor-input constraint; no-history and no-anchor variants identify what is lost. & C1--C3 \\
E8 & S9--S10: real indoor station-keeping and unsafe-command tests & B0, B3, B5, Ours & mocap/evaluation-only drift, range-anchor error, min clearance, intervention rate, repeated-trial success & Hardware trends match simulation: observable geometries allow lower drift; degenerate scenes drift predictably; unsafe commands are attenuated before violating the safety boundary. & C1--C4 \\
\bottomrule
\end{tabular}
\end{table*}

%============================================================
\subsection{Expected Results by Contribution}
\label{sec:exp_expected_by_contribution}
%============================================================

\subsubsection{C1: Pose-, odometry-, and velocity-free learned governor}
The main evidence for C1 is not that the method outperforms an odometry-fed
upper bound.  Rather, the evidence is that it remains functional without those
inputs.  E3 and E7 should show that the proposed governor achieves stable
command tracking and station-keeping under the audited input constraint, while
B1/B2 quantify the performance gap when explicit state information is allowed.
A small gap indicates that range--inertial history captures much of the
short-horizon motion information needed for local regulation; a larger gap is
acceptable if the method still satisfies the station-keeping and safety trends
claimed in C2--C4.

\subsubsection{C2: Measurement-space station-keeping}
E1 should show that the range anchor reduces both physical drift and range-anchor
error in scenes where the local geometry provides observable range variation.
B3 is the critical ablation: if removing the anchor does not increase drift or
range error, then the proposed station-keeping mechanism is not supported.  B0
separates the effect of the low-level velocity controller from the learned
measurement-space governor.

\subsubsection{C3: Observability-limited holding}
E2 should show that station-keeping performance follows the conditioning of the
range Jacobian.  The desired outcome is not perfect holding in every scene, but
a consistent ordering: corners should produce the smallest drift, flat walls and
corridors should exhibit drift in weakly observable directions, and open-field
no-return scenes should show the largest drift.  This experiment is the key
evidence that station-keeping failure is a predictable geometric limitation
rather than merely poor training.

\subsubsection{C4: ICS-inspired command intervention}
E4--E6 evaluate whether intervention occurs when a command exceeds local
braking feasibility and remains inactive when the command is feasible.  Compared
with B5, the full method should reduce collisions and ICS-violation time.
Compared with B6, it should avoid late distance-only reactions.  Compared with
B7, it should preserve tracking better by using learned context rather than a
fixed attenuation rule alone.

%============================================================
\subsection{Real-Robot Validation Protocol}
\label{sec:exp_real_robot}
%============================================================

The real-robot experiments are deliberately limited to cases that directly
support the paper's claims.  Each trial is repeated with the same initial
command script and with randomized initial yaw and lateral offset within a safe
range.  Motion capture, external tracking, or a safety pilot may be used for
logging and emergency supervision, but the deployed governor receives only the
same range--inertial history and command inputs used in simulation.

The minimum real-robot protocol consists of three tests.
\begin{enumerate}
    \item \textbf{Observable station-keeping:} the vehicle is initialized near a
    corner or inside a short corridor, receives a null command for $30$--$60$ s,
    and the drift of Ours is compared with B0 and B3.
    \item \textbf{Degenerate station-keeping:} the vehicle is placed near a flat
    wall or in a corridor-aligned configuration.  The expected result is not zero
    drift, but drift concentrated along the weakly observable direction.
    \item \textbf{Unsafe command attenuation:} the operator command points toward
    a wall at increasing commanded speeds.  The test reports command attenuation,
    minimum clearance, and intervention timing.  The trial is terminated by a
    safety supervisor if the vehicle reaches a predefined hard safety boundary.
\end{enumerate}

A real-robot result is considered consistent with the paper if the qualitative
ordering observed in simulation is reproduced: measurement-space anchoring
improves holding in observable geometries; degeneracy produces predictable
rather than random drift; and the ICS-inspired layer attenuates commands before
minimum-clearance violation.  The physical experiments are not used to claim
formal safety guarantees, since the deployed system still relies on the
low-level controller, sensing latency, and hardware-specific emergency stops.

%============================================================
\subsection{Reporting and Statistical Treatment}
\label{sec:exp_reporting}
%============================================================

For each simulated setting we report the mean, standard deviation, and
interquartile range over at least $N_{\mathrm{sim}}$ randomized episodes.  For
hardware experiments we report all trials, including failures and emergency
stops, and use median and range when the number of trials is small.  Collision
rates are reported with binomial confidence intervals.  Drift and tracking
metrics are computed over identical command scripts for all compared methods.

To avoid overstating the pose-free claim, every table includes an input column
indicating whether the method receives pose, odometry, explicit velocity,
reconstructed maps, or only range--inertial history.  Evaluation-only ground
truth is reported separately and is never counted as a policy input.


%============================================================
\section{Conclusion}
\label{sec:conclusion}
%============================================================

This paper presented a pose-free command governor for quadrotor
station-keeping and safe velocity tracking from range--inertial history.  The
formulation treats null-command holding and nonzero-command tracking as a
single measurement-space regulation problem, clarifies when such regulation is
geometrically observable, and attenuates commands that violate local braking
feasibility.  The resulting protocol separates what the deployed actor may use
from what the simulator or evaluator may measure, making drift, degeneracy, and
safety interventions explicit rather than hidden behind a pose estimator.

%============================================================
% BIBLIOGRAPHY  (thebibliography with \bibitem)
%============================================================
\begin{thebibliography}{1}

\bibitem{sprind2025}
Authors,
``Kilometer-scale GNSS-denied UAV navigation via heightmap gradients:
a winning system from the SPRIN-D challenge,''
\emph{arXiv preprint arXiv:2510.01348}, 2025.

\bibitem{neuralcbf2024}
Authors,
``Neural control barrier functions for safe navigation,''
\emph{arXiv preprint arXiv:2407.19907}, 2024.

\bibitem{locus2023}
Authors,
``LiDAR-based online control barrier function synthesis for safe navigation
in unknown environments,''
in \emph{Proc. IEEE/RSJ Int. Conf. Intelligent Robots and Systems (IROS)},
2023.

\bibitem{compositecbf2025}
Authors,
``Safe quadrotor navigation using composite control barrier functions,''
\emph{arXiv preprint arXiv:2502.04101}, 2025.

\bibitem{reactive2026}
Authors,
``Reactive robot-centric safety for autonomous navigation in constrained and
dynamic environments,''
\emph{arXiv preprint arXiv:2605.15782}, 2026.

\end{thebibliography}

\end{document}